\section{Complete Computational Algorithm}

\begin{algorithm}[H]
\caption{End-to-end hybrid tsunami reconstruction and future calibration workflow.}
\label{alg:complete-workflow}
\begin{algorithmic}[1]
\State Establish source provenance governance and approved observation registry.
\State Ingest candidate offshore, nearshore, run-up and inundation observations after approval.
\State Define numerical convergence plan for Regional2D and Local3D mesh/timestep controls.
\State Acquire and verify terrain, bathymetry and finite-fault source artefacts.
\State Transform event and corridor geometry to EPSG:32654.
\State Generate earthquake displacement \(\Delta z_b\) and initialise Regional2D free surface.
\State Solve the Regional2D finite-volume model with wet/dry, friction and boundary treatments.
\State Extract virtual gauges and nearshore coupling histories.
\State Perform source-model sensitivity after convergence and observation approval.
\State Calibrate Regional2D physical parameters only through documented hypotheses.
\State Execute Regional holdout validation without arbitrary time shifts.
\State Select accepted coupling window and project normal/tangential momentum.
\State Reconstruct 3D inlet velocity and VOF boundary data.
\State Generate Local3D geometry, mesh, OpenFOAM dictionaries and provenance records.
\State Run URANS--VOF no-defence and defence cases.
\State Extract free surface, probes, forces and moments.
\State Perform Local3D turbulence, roughness and geometry sensitivity as authorised.
\State Calibrate Local3D physical parameters only after Regional forcing is accepted.
\State Execute hybrid holdout validation.
\State Archive HDF5/VTKHDF results, publication figures and provenance.
\State Record retain/revert decisions and update the evidence ledger.
\end{algorithmic}
\end{algorithm}

The G6 work completed the implementation and real-event integration path through Local3D acceptance. The remaining algorithmic steps define validation, calibration, archive and publication work.
